% !PRINTED 550
\weberpage{550}
%<<CONT
\weq{17}{\sqrt[m]{\psi_{\lambda}} = \chi_{\lambda}\left(\sqrt{a_1^{m_1} a_2^{m_2} \ldots}\right)^{-\nu} e^{-\frac{i\Theta}{m}\nu}.}

By (15) we can set
\[
e^{i\Theta_1} = \frac{b_1 + ic_1}{\sqrt{a_1^m}}, \qquad e^{i\Theta_2} = \frac{b_2 + ic_2}{\sqrt{a_2^m}} \cdots,
\]
and if, then, by decomposition into the real and imaginary part we
obtain
\weq{18}{(b_1 + ic_1)^{m_1}(b_2 + ic_2)^{m_2} \cdots = B + iC}
and set
\weq{19}{a_1^{m_1} a_2^{m_2} \cdots = A}
it follows by (16) that
\[
e^{i\Theta} = \frac{B + iC}{\sqrt{A^m}},
\]
hence
\weq{20}{\sqrt{A^m}\cos\Theta = B, \qquad \sqrt{A^m}\sin\Theta = C,}
and
\weq{21}{\sqrt[m]{\psi_{\lambda}} = \chi_{\lambda}\sqrt{A}^{-\nu}\left(\cos\frac{\Theta\nu}{m} - i\sin\frac{\Theta\nu}{m}\right).}

The angle $\Theta$ is determined by (20) only up to a multiple of
$2\pi$, and if one therefore puts $\Theta + 2h\pi$ for $\Theta$ and lets $h$
run from $0$ to $m - 1$, one obtains from (21) the $m$ different values of the root quantity. The functions $\cos\dfrac{\Theta\nu}{m}$
and $\sin\dfrac{\Theta\nu}{m}$ can further be expressed rationally in terms of $\cos\dfrac{\Theta}{m}$, $\sin\dfrac{\Theta}{m}$. The solution of the cyclic equations in
the real field $\Omega$ is thus reduced to the following.

\gesperrt{One first adjoins to the field $\Omega$ the $m$th root of unity $\varepsilon$ (division of the circumference of the circle into $m$ equal
parts). Thereupon $A, B, C$ are known. One further adjoins
the positive square root $\sqrt{A}$; then
$\cos\Theta$ and $\sin\Theta$ are known by (20). Finally one adjoins
$\cos\dfrac{\Theta}{m}$, $\sin\dfrac{\Theta}{m}$ (division of the angle $\Theta$ into $m$ parts).
Then the cyclic equation is solved by (21).}

These considerations lead further to an interesting
result on the reality relations of the roots of cyclic
equations.
